
%(BEGIN_QUESTION)
% Copyright 2006, Tony R. Kuphaldt, released under the Creative Commons Attribution License (v 1.0)
% This means you may do almost anything with this work of mine, so long as you give me proper credit

Reguleringsventiltyper (f.eks. globe, kule, spjeld) kan gis relative kapasitetsverdier kalt $C_d$-faktor. Dette ligner konseptet med strømningskapasitet ($C_v$), men er generalisert for en bestemt type eller konstruksjon av reguleringsventil. Ligningen som knytter $C_d$ til $C_v$ er:

$$C_d = {C_v \over d^2}$$

\noindent
Der,

$C_d$ = Relativ strømningskapasitet for ventiltypen

$C_v$ = Maksimal strømningskapasitet for den aktuelle ventilen

$d$ = Nominell rørdimensjon for den aktuelle ventilen, i tommer

\vskip 10pt

Flere ventilkapasitetsfaktorer ($C_d$) for ulike reguleringsventiltyper er vist her, forutsatt full trimareal og helt åpen posisjon:

% Ingen blanklinjer er tillatt mellom linjer i en \halign-struktur!
% Jeg bruker kommentarer (%) i stedet, slik at TeX ikke feiler.

$$\vbox{\offinterlineskip
\halign{\strut
\vrule \quad\hfil # \ \hfil &
\vrule \quad\hfil # \ \hfil \vrule \cr
\noalign{\hrule}
%
% Første rad
{\bf Ventilkonstruksjon} & $C_d$ \cr
%
\noalign{\hrule}
%
% Neste rad
Enkeltsetet globeventil, portet plugg & 9.5 \cr
%
\noalign{\hrule}
%
% Neste rad
Enkeltsetet globeventil, profilert plugg & 11 \cr
%
\noalign{\hrule}
%
% Neste rad
Enkeltsetet globeventil, karakteriserende bur & 15 \cr
%
\noalign{\hrule}
%
% Neste rad
Dobbeltsetet globeventil, portet plugg & 12.5 \cr
%
\noalign{\hrule}
%
% Neste rad
Dobbeltsetet globeventil, profilert plugg & 13 \cr
%
\noalign{\hrule}
%
% Neste rad
Roterende kuleventil, segmentert & 25 \cr
%
\noalign{\hrule}
%
% Neste rad
Roterende kuleventil, standard port (diameter $\approx$ 0.8$d$) & 30 \cr
%
\noalign{\hrule}
%
% Neste rad
Roterende spjeldventil, 60$^{o}$, uten forskjøvet sete & 17.5 \cr
%
\noalign{\hrule}
%
% Neste rad
Roterende spjeldventil, 90$^{o}$, med forskjøvet sete & 29 \cr
%
\noalign{\hrule}
%
% Neste rad
Roterende spjeldventil, 90$^{o}$, uten forskjøvet sete & 40 \cr
%
\noalign{\hrule}
} % Slutt på \halign
}$$ % Slutt på \vbox

Disse $C_d$-dataene for ulike reguleringsventiltyper gjør det mulig å tilnærme en ventils fullstrøms-$C_v$ når vi kjenner ventiltype og rørdimensjon. Beregn fullstrøms-$C_v$ for disse reguleringsventilene:

\begin{itemize}
\item{} Segmentert kuleventil, 4 tommers rørdimensjon; $C_v$ = \underbar{\hskip 50pt}
\vskip 10pt
\item{} Enkeltsetet, burført globeventil, 6 tommers rørdimensjon; $C_v$ = \underbar{\hskip 50pt}
\vskip 10pt
\item{} Dobbeltsetet, portet-plugg globeventil, 2 tommers rørdimensjon; $C_v$ = \underbar{\hskip 50pt}
\vskip 10pt
\item{} 90$^{o}$ spjeldventil med forskjøvet sete, 20 tommers rørdimensjon; $C_v$ = \underbar{\hskip 50pt}
\end{itemize}

\vskip 20pt \vbox{\hrule \hbox{\strut \vrule{} {\bf Forslag til sokratisk diskusjon} \vrule} \hrule}

\begin{itemize}
\item{} Gitt en bestemt nødvendig $C_v$-verdi: Hvilken reguleringsventilkonstruksjon tillater minst ventilstørrelse (rørdimensjon)?
\item{} Gitt en bestemt nødvendig $C_v$-verdi: Hvilken reguleringsventilkonstruksjon krever størst ventilstørrelse (rørdimensjon)?
\end{itemize}

\underbar{file i01371}
%(END_QUESTION)





%(BEGIN_ANSWER)

\begin{itemize}
\item{} Segmentert kuleventil, 4 tommers rørdimensjon; $C_v$ = \underbar{\bf 400}
\vskip 5pt
\item{} Enkeltsetet, burført globeventil, 6 tommers rørdimensjon; $C_v$ = \underbar{\bf 540}
\vskip 5pt
\item{} Dobbeltsetet, portet-plugg globeventil, 2 tommers rørdimensjon; $C_v$ = \underbar{\bf 50}
\vskip 5pt
\item{} 90$^{o}$ spjeldventil med forskjøvet sete, 20 tommers rørdimensjon; $C_v$ = \underbar{\bf 11,600}
\end{itemize}

\vskip 10pt

Merk: $C_v$-verdier beregnet ved bruk av relative strømningskoeffisienter ($C_d$) er {\it kun tilnærminger}! Denne beregningsmetoden bør bare brukes til å {\it estimere} nødvendig ventilstørrelse for en bestemt applikasjon.

%(END_ANSWER)





%(BEGIN_NOTES)

Ligningen studentene må bruke for å finne $C_v$ i hvert tilfelle er selvsagt:

$$C_v = d^2 C_d$$

\vskip 10pt

Kilde for $C_d$-faktorer: \underbar{Chapter 4.17: Valve Sizing}; {\it Instrument Engineer's Handbook, Process Control, Third Edition}; side 590.

%INDEX% Final Control Elements, valve: relative flow capacity (defined)
%INDEX% Final Control Elements, valve: relative flow coefficient (defined)
%INDEX% Final Control Elements, valve: sizing

%(END_NOTES)

